% 한글 초안 (수업 과제용) -> 이후 영어로 옮겨 RSS 제출.
% Overleaf: Menu > Compiler = XeLaTeX (kotex 필요).
% \todo{...} 는 아직 채우지 않은 부분. 최종본 전에 모두 지운다.
% 수치는 저장소의 실험 기록에서 그대로 옮긴 것. 출처는 각 줄 옆 주석(%)에 적었다.
% 2026-11-26 방향 변경: "다음 작업 확정" 제안 -> "기록 불일치 창고에서의 학습 배정", 확정은 부정적 제거 실험으로.
\documentclass[10pt,twocolumn]{article}
\usepackage{kotex}
\usepackage[margin=0.75in]{geometry}
\usepackage{amsmath,amssymb}
\usepackage{booktabs}
\usepackage{graphicx}
\usepackage{xcolor}
\usepackage{hyperref}
\newcommand{\todo}[1]{\textcolor{red}{[TODO: #1]}}

\title{재고 기록이 어긋나는 물류창고에서의\\경량 시뮬레이터 기반 다중 로봇 작업 배정 학습}
\author{(이중 블라인드: 저자 정보 생략)}
\date{}

\begin{document}
\maketitle

% ---------------------------------------------------------------
\begin{abstract}
물류창고의 창고 실행 시스템(WES)은 입고 적치, 출고 피킹, 재고 확인 작업을 여러 이동 로봇에게 실시간으로 배정한다.
현장 작업자가 기록 없이 재고를 옮기거나 빼 가면 시스템 기록과 실제 재고가 어긋나고, 로봇은 빈 칸에 가거나 확인 작업을 추가로 해야 한다.
본 연구는 작업-로봇 짝에 선형 점수를 매기는 배정 정책의 가중치를, 실제 WMS·WES 코드를 그대로 쓰고 로봇 이동만 단순화한 경량 시뮬레이터에서 교차 엔트로피 방법(CEM)으로 학습하고, 물리 시뮬레이터 Webots에서 검증한다.
실제 신발 제조사 창고의 주문 데이터로 만든 시나리오에서, 작업자 오류가 있을 때 학습한 배정은 팀 기본 WES보다 주문 리드타임을 경량 시뮬레이터에서 7.4\% 줄였다. 이는 오류로 생긴 손실보다 크다.
Webots의 오류 조건 18개 시험 시나리오에서 학습한 배정은 주문 리드타임을 11--22\% 줄였다. 다만 개선 폭은 학습된 가중치에 따라 달랐고, 과부하 조건에서는 한 가중치가 큰 주문을 끝내지 못하는 경우가 있었다.
또한 일하고 있는 로봇에게 다음 작업을 미리 예약하는 \emph{다음 작업 확정}은 성능을 높이지 못했으며, 학습은 이 선택지를 쓰지 않는 쪽으로 수렴했다.
\end{abstract}

% ---------------------------------------------------------------
\section{서론}

자동화 물류창고에서는 여러 대의 이동 로봇이 입고된 상자를 랙에 적치하고(putaway), 주문에 맞춰 상자를 출고 레인으로 옮기며(pick),
기록과 실제가 다른지 의심되는 칸을 확인한다(verify)~\cite{wurman2008}.
어떤 작업을 어느 로봇에게 언제 맡길지 정하는 작업 배정은 창고의 처리 속도와 출고 마감 준수에 직접 영향을 준다~\cite{gerkey2004}.

실제 창고에서는 로봇과 사람이 함께 일한다.
작업자가 상자를 기록 없이 다른 칸으로 옮기거나 빼 가면 재고 기록이 틀어지고~\cite{dehoratius2008},
로봇은 상자가 없는 칸에 도착해 피킹에 실패하거나, 이미 차 있는 칸에 적치하려다 다른 칸을 찾아야 한다.
그 결과 확인 작업과 재작업이 늘어나고, 같은 배정 규칙이라도 성능이 떨어진다.
본 연구의 경량 시뮬레이터에서 작업자 오류를 넣으면 팀 기본 WES의 주문 리드타임은 495\,초에서 527\,초로 늘었다. % logs/train_nofault_m100_report.json, logs/train_fault_base_cem_r0_report.json stand-in

배정 규칙은 보통 사람이 정한 우선순위(가까운 로봇, 피킹 우선 등)로 만든다.
이런 규칙은 기록이 정확하다는 가정에서 조정되었기 때문에, 확인 작업과 재작업이 섞이는 상황에서 어떤 작업을 먼저 할지를 잘 고르지 못한다.
본 연구는 배정 점수의 가중치를 데이터 기반 시나리오에서 직접 학습하여 이 문제를 다룬다.

\textbf{기여.}
\begin{enumerate}
  \item 실제 창고 주문 데이터와 작업자 오류 모델로 만든 시나리오에서, 실제 WMS·WES 코드를 그대로 쓰는 경량 시뮬레이터로 배정 가중치를 학습하고 물리 시뮬레이터(Webots)에서 검증하는 절차를 제시한다.
  \item 작업자 오류가 있을 때 학습한 배정이 팀 기본 WES보다 주문 리드타임과 p90을 줄임을 경량 시뮬레이터(학습 시드 3개, 두 창고 규모)와 Webots(18개 시나리오)에서 보인다. 동시에 학습된 가중치 사이의 차이, 과부하에서의 실패, 거리 최적 매칭(헝가리안)이라는 강한 비교군을 함께 보고한다.
  \item 일하고 있는 로봇에게 다음 작업을 예약하는 \emph{다음 작업 확정}을 제거 실험으로 분석한다. 강제 확정은 성능을 낮추고, 학습은 확정을 거의 쓰지 않는 가중치로 수렴한다. 탐색 방법(CMA-ES)이나 도착 예상 시간 특징을 바꿔도 결과는 비슷했다. 즉 이 문제에서 개선은 \emph{어느 로봇}이 아니라 \emph{어떤 작업을 먼저} 할지에서 나온다.
\end{enumerate}

% ---------------------------------------------------------------
\section{관련 연구}

\textbf{다중 로봇 작업 배정.}
Gerkey와 Matari\'{c}~\cite{gerkey2004}은 다중 로봇 작업 배정 문제를 분류했다.
Choudhury 등~\cite{choudhury2020}은 완료 시간의 불확실성과 시간 제약이 있는 동적 배정을 다뤘고,
Dhanaraj 등~\cite{dhanaraj2024}은 실패와 복구 작업을 고려한 hindsight optimization을 제안했다.
\todo{우리 방법과의 차이: 이들은 작업 실패 확률을 알고 계획하지만, 우리는 기록 불일치가 언제 어디서 생기는지 모르는 상태에서 WMS가 보여 주는 정보만으로 배정한다}

\textbf{물류창고 로봇.}
Wurman 등~\cite{wurman2008}은 수백 대 규모 창고 로봇 시스템을 소개했고,
Leet 등~\cite{leet2023}은 작업 구성, 배정, 스케줄, 경로 계획을 결합해 대규모 작업량을 다뤘다.
본 연구는 그보다 작은 규모에서, 기록 불일치가 있는 조건의 배정 정책 학습에 집중한다.

\textbf{재고 기록 불일치.}
실제 재고와 시스템 기록이 다른 문제는 소매·물류에서 널리 보고되었다~\cite{dehoratius2008}.
검사와 보충을 결합한 정책~\cite{kok2007,akkerman2025}이 연구되었으나, 로봇 이동과 자원 경합을 함께 다루지는 않는다.

\textbf{정책 탐색.}
교차 엔트로피 방법(CEM)~\cite{rubinstein1999}과 CMA-ES~\cite{hansen2001}는 경사 없이 정책 파라미터를 찾는 방법이다.
에피소드 끝에 한 번 나오는 비용과 적은 파라미터 수에 적합하다.

% ---------------------------------------------------------------
\section{문제 설정}

\subsection{창고와 작업}
기본 창고는 $24\times14$\,m 크기이며 랙 칸, 입고 대기칸(buffer), 지역별 출고 레인 4개(각 8칸)로 구성된다. 로봇은 8대이다. % scenario.json world, outbound.lanes
규모 확장 실험에는 $34\times20$\,m, 랙 480칸, 로봇 14대의 큰 창고를 쓴다. % gen_world --layout large --transports 14
작업 종류는 다음과 같다.
\begin{itemize}
  \item \textbf{PUTAWAY}: 입고 대기칸의 상자를 정해진 랙 칸에 적치한다.
  \item \textbf{PICK}: 주문 품목이 있는 랙 칸의 상자를 해당 지역 출고 레인으로 옮긴다.
  \item \textbf{RELOCATE}: 상자를 다른 칸으로 옮긴다. 배정에서는 적치와 같은 종류로 다룬다.
  \item \textbf{VERIFY}: 칸의 실제 내용을 확인해 기록을 바로잡는다. 놀고 있는 로봇이 수행한다.
\end{itemize}
출고 트럭은 지역별로 900\,초마다 마감되어 출발한다. % cutoff_period_s
에피소드 길이는 3600\,초이다.

\subsection{시스템 구조}
창고 관리 시스템(WMS)은 재고 기록과 작업 상태를 관리하고 작업을 만든다.
WES는 WMS가 보여 주는 정보만으로 매 단계 작업과 로봇을 짝지어 실행 지시를 보낸다.
WES는 실제 재고나 미래 주문 같은 숨겨진 정보를 쓰지 않는다.
\todo{그림 1: WMS--WES--로봇 구조도}

\subsection{시나리오}
주문은 공개된 신발 제조사 창고 피킹 데이터셋~\cite{assis2023}의 \texttt{Customer\_Order.csv}에서 하루를 골라 시간을 압축해 만든다.
\begin{itemize}
  \item 시간 압축 비율은 약 6배이고, 웨이브는 90\,초 간격으로 풀린다. % time_scale 6.08, wave_interval_s 90
  \item 입고는 그날 수요를 덮도록 합성한 트럭으로 만들고, 트럭 한 대에 상자 12개가 실린다.
  \item 부하는 2023-03-09(보통 물량), 2023-03-22의 주문 50\% 표본(절반 부하), 2023-03-22 전체(과부하)의 세 가지다.
\end{itemize}

\subsection{평가 지표}
\begin{itemize}
  \item \textbf{주문 리드타임}(주 지표): 주문이 풀린 시각부터 그 주문의 마지막 상자가 출하된 시각까지의 평균이다. 끝내지 못한 주문은 에피소드 종료 시각을 완료 시각으로 쳐서 포함한다. 끝낸 주문만 평균하면 어려운 주문을 미루는 정책이 유리해 보이기 때문이다.
  \item \textbf{상자 단위 리드타임}: 같은 방식으로 상자마다 잰 평균이다.
  \item \textbf{p90}: 주문 리드타임의 90번째 백분위수이다.
  \item \textbf{출하 수}와 \textbf{마감 놓침}(마감까지 출하하지 못한 주문 수)이다.
\end{itemize}

\subsection{사람 작업자에 의한 기록 불일치}
현장 작업자가 기록 없이 재고를 옮기는 상황을 세 가지로 모델링한다.
\begin{itemize}
  \item S1: 상자를 다른 칸으로 옮긴다. 시간당 8회, 한 번에 최대 3칸이다.
  \item S2: 상자를 빼 간다. 시간당 3회이다.
  \item S3: 적치하려는 칸이 이미 차 있다. 확률 0.1이다.
\end{itemize}
불일치는 로봇이 칸에 도착해 라벨을 읽을 때 드러난다.
피킹 실패(SHORT\_PICK)나 품목 불일치(ITEM\_MISMATCH)가 나면 이웃 칸을, 적치 칸이 차 있으면(DESTINATION\_OCCUPIED) 그 칸을 확인 작업으로 만든다(열린 확인 작업 최대 4개). % sched_verify VERIFY_POLICY=event, VERIFY_MAX_OPEN 4
로봇 자체의 오류(적치 실수, 스캐너 실패)는 0으로 둔다.
실험은 오류 없음과 위 기본값의 두 단계로 한다.
위 비율은 실측 자료가 없어 정한 합성 가정이며, 한계에서 다시 언급한다.

% ---------------------------------------------------------------
\section{방법}

\subsection{작업-로봇 짝 점수}
매 단계 시작 가능한 대기 작업 $\tau$ 각각을, 놀고 있는 로봇과 일하고 있는 로봇 $r$ 각각과 짝지어 점수를 매긴다.
\begin{equation}
  S(\tau, r) = \mathbf{w}^\top \boldsymbol{\phi}(\tau, r)
\end{equation}
점수가 높은 짝부터 탐욕적으로 고르며, 한 작업과 한 로봇은 한 번만 쓰인다.
놀고 있는 로봇이 이기면 바로 실행 지시를 보내고, 일하고 있는 로봇이 이기면 아무 지시도 보내지 않고 다음 단계에서 다시 정한다.
특징 $\boldsymbol{\phi}$는 표~\ref{tab:features}와 같다.
일하고 있는 로봇의 거리는 \emph{현재 작업의 남은 이동 거리 + 현재 작업이 끝나는 곳에서 새 작업까지의 거리}로 계산한다.
모든 특징은 WMS가 보여 주는 정보로만 계산한다.

\begin{table}[t]
\centering
\caption{배정 점수의 특징. 위 세 개만 로봇에 따라 달라진다.} % lib/sched_learned.py FEATURES
\label{tab:features}
\small
\begin{tabular}{ll}
\toprule
특징 & 정의 \\
\midrule
dist & 로봇에서 작업 출발 칸까지 거리 (m) / 10 \\
busy & 일하고 있는 로봇이면 1 \\
eta & 로봇이 작업에 도착할 예상 시간 (분) \\
\midrule
pick, putaway, verify & 작업 종류 (0/1) \\
age & 대기 시간 / 600\,s (최대 2) \\
slack & 마감까지 남은 시간 / 900\,s, [0,1] (PICK) \\
truck\_here & 해당 레인에 트럭이 있으면 1 (PICK) \\
lane\_fill & 출고 레인 점유율 (PICK) \\
buffer\_fill & 입고 대기칸 점유율 (PUTAWAY) \\
\bottomrule
\end{tabular}
\end{table}

\subsection{경량 시뮬레이터}
Webots 실행은 1시간 시나리오에 약 8.5분이 걸려 학습에 쓰기 어렵다.
그래서 같은 시나리오 파일과 같은 WMS·WES 코드를 쓰고, 로봇 이동만 단순화한 경량 시뮬레이터를 만들었다.
작업자 오류도 같은 코드로 들어가므로, 학습 중에 기록 불일치와 확인 작업이 그대로 나타난다.
\begin{itemize}
  \item 한 구간의 이동 시간은 \emph{경로 길이 / 속도 + 2\,초 + 1\,초 $\times$ (우선권이 앞선 움직이는 로봇 수)}이다.
  \item 시간 간격은 0.5\,초이고, 충돌과 회피는 모델링하지 않는다.
  \item 1시간 시나리오를 약 1초에 실행한다.
\end{itemize}
Webots 실행 51개와 비교했을 때 가동률 평균 오차는 2.2\%, 주문 리드타임 평균 오차는 5.4\%였다. % docs/lite_sim.md, 2026-11-16
방법 간 순위는 리드타임 기준 92\%, 출하 기준 86\%가 일치했다.
\todo{이 보정은 Webots에 비켜서기 동작을 넣기 전 결과. 재보정하거나 한계에 명시}

\subsection{학습}
가중치 $\mathbf{w}$는 경량 시뮬레이터에서 CEM으로 학습한다.
\begin{itemize}
  \item 후보 24개, 상위 6개, 15회 반복, 초기 표준편차 1.0(최소 0.05)이다.
  \item 시작값은 \emph{가까운 로봇, 피킹 우선} 규칙이다(dist $-1$, busy $-1$, pick $+2$).
  \item 매 반복 모든 후보를 학습 시나리오 12개 전부에서 평가해 같은 조건으로 비교한다.
  \item 학습 난수 시드를 3개(r0--r2) 바꿔 결과의 흔들림을 잰다.
\end{itemize}
에피소드 비용은 다음과 같다.
\begin{equation}
  C = 0.5\,L_{\text{order}} + 0.5\,L_{\text{unit}} + \lambda\,N_{\text{missed}}, \quad \lambda = 100
\end{equation}
$\lambda=200, 300$은 시험 시나리오에서 기본 방식보다 나빠져 제외했다.

\subsection{제거 실험: 다음 작업 확정}
위 구조에서는 일하고 있는 로봇이 짝을 이겨도 지시가 나가지 않으므로, 로봇이 일을 마치는 순간에 다시 결정된다.
곧 일을 마치는 가까운 로봇에게 다음 작업을 맡기는 효과를 보기 위해 \emph{다음 작업 확정}을 넣어 비교한다.
일하고 있는 로봇 $r$이 작업 $\tau$를 이기면 $\tau$를 $r$의 다음 작업으로 예약한다.
\begin{itemize}
  \item 예약된 작업은 다른 로봇의 후보에서 빠지고, $r$이 일을 마치면 예약된 작업을 실행한다.
  \item 작업이 대기 상태를 벗어나거나 시작할 수 없으면 예약은 풀린다. 로봇 한 대에는 예약이 하나만 걸린다.
  \item 두 가지 변형을 쓴다. \textbf{항상 확정}은 일하는 로봇이 이기면 늘 예약한다. \textbf{조건부 확정}은 같은 역할의 놀고 있는 로봇이 남아 있는데도 일하는 로봇이 이겼을 때만 예약한다.
\end{itemize}
예약 여부는 같은 점수 $S$로 정해지므로, 학습은 확정을 쓸지 말지도 가중치로 배운다.

% ---------------------------------------------------------------
\section{실험 설정}

\textbf{시나리오.}
\begin{itemize}
  \item 학습: 03-09와 03-22 절반, 각각 시드 42--47 (12개)
  \item 시험(경량): 03-09와 03-22 절반은 시드 48--50, 03-22 전체(학습에 없던 과부하)는 시드 42--44 (9개), 시나리오마다 5회 실행 (45 에피소드)
  \item 오류 조건과 큰 창고는 같은 날짜·시드로 따로 만든다.
\end{itemize}

\textbf{비교 방법.}
\begin{itemize}
  \item 팀 기본 WES 배정
  \item 가까운 로봇, 피킹 우선 규칙 (학습 시작값)
  \item 헝가리안 최적 매칭~\cite{kuhn1955}: 매 단계 놀고 있는 로봇과 시작 가능한 작업(같은 후보 30개)을, 로봇에서 작업 출발 칸까지 거리의 합이 최소가 되도록 짝짓는다. 작업 우선순위는 쓰지 않는다. % lib/sched_hungarian.py
  \item 제안 방법: CEM 학습 배정
  \item 제거 실험: 확정 + CEM, CMA-ES, 도착 예상 시간 특징
\end{itemize}

\textbf{Webots 검증.}
로봇 8대, 바퀴 물리와 충돌 회피, 통행 우선순위 규칙을 포함한다.
작업자 오류 기본값에서 세 부하 $\times$ 학습에 쓰지 않은 시드 48--50 (9개 시나리오)을 각 방법으로 한 번씩 실행한다. % scratchpad webots_fault.sh, 2026-11-26
비교 방법은 팀 기본 WES, 헝가리안, 오류 조건에서 학습한 CEM(r0), 오류 없이 학습한 CEM이다.
2차 실행에서 오류 조건 시드 51--53과 헝가리안, 오류 없음 시드 48--50을 더했다. % scratchpad webots_more.sh, 2026-11-26

% ---------------------------------------------------------------
\section{결과}

\subsection{작업자 오류 조건 (경량 시뮬레이터)}

\begin{table}[t]
\centering
\caption{작업자 오류 기본값, 경량 시뮬레이터 시험 결과 (45 에피소드 평균). 학습 결과는 학습 시드 3개의 평균 $\pm$ 표준편차. 시간은 초.} % logs/train_fault_{base,large}_cem_r{0,1,2}_report.json
\label{tab:fault}
\small
\begin{tabular}{lrrrr}
\toprule
방법 & 주문 & 상자 & p90 & 마감 놓침 \\
\midrule
\multicolumn{5}{l}{\emph{기본 창고 (로봇 8대)}} \\
팀 기본 WES & 527 & 423 & 753 & 2.29 \\
가까운 로봇 규칙 & 508 & 416 & 718 & 2.27 \\
헝가리안 & 505 & 415 & 696 & 2.4 \\ % scratchpad hung_screen.py, 2026-11-26
CEM & 488$\pm$8 & 418$\pm$2 & 659$\pm$29 & 2.2 \\
\midrule
\multicolumn{5}{l}{\emph{큰 창고 (로봇 14대)}} \\
팀 기본 WES & 508 & 422 & 697 & 3.0 \\
가까운 로봇 규칙 & 498 & 416 & 701 & 2.6 \\
CEM & 477$\pm$4 & 416$\pm$2 & 621$\pm$36 & 2.7 \\
\bottomrule
\end{tabular}
\end{table}

\begin{itemize}
  \item 기본 창고에서 CEM 배정은 같은 에피소드끼리 비교했을 때 팀 기본 WES보다 주문 리드타임이 평균 39\,초(7.4\%) 짧았다(학습 시드 간 표준편차 8\,초).
  \item 작업자 오류로 팀 기본 WES는 32\,초(495$\to$527) 나빠졌다. CEM의 개선 폭은 이 손실보다 크다.
  \item 오류 없는 조건에서 CEM의 개선은 22\,초였으므로(표~\ref{tab:nofault}), 경량 시뮬레이터에서는 기록 불일치가 있을 때 학습의 이득이 더 컸다. 이 차이는 Webots에서 확인되지 않았다(\S\ref{sec:webots}).
  \item p90은 753에서 659\,초로 줄었다. 출하 수와 마감 놓침은 거의 같았다.
  \item 큰 창고에서도 31\,초 개선이 유지되었다.
  \item 헝가리안 최적 매칭은 경량 시뮬레이터에서 팀 기본 WES보다 23\,초 짧은 데 그쳤고 오류 없는 조건에서는 차이가 없었다($+1$\,초). Webots에서는 이보다 강했다(\S\ref{sec:webots}).
\end{itemize}

\subsection{Webots 검증}
\label{sec:webots}
\begin{table*}[t]
\centering
\caption{Webots 결과. 시간은 초, 시나리오마다 1회 실행. 부하별 값은 주문 리드타임의 시드 평균. 마지막 열은 같은 시나리오끼리 잰 주문 리드타임 차이의 평균 $\pm$ 표준오차와, 팀 기본 WES보다 나았던 시나리오 수.} % logs/*_fl_*, *_nf2_*; scratchpad agg2.py, 2026-11-27
\label{tab:webots}
\small
\begin{tabular}{lrrrrrrrr}
\toprule
 & \multicolumn{3}{c}{부하별 주문 리드타임} & \multicolumn{4}{c}{전체} & \\
\cmidrule(lr){2-4}\cmidrule(lr){5-8}
방법 & 03-09 & 03-22 절반 & 03-22 전체 & 주문 & 상자 & p90 & 마감 놓침 & 팀 WES 대비 \\
\midrule
\multicolumn{9}{l}{\emph{작업자 오류 기본값, 시드 48--53 (18개)}} \\
팀 기본 WES & 549 & 525 & 467 & 514 & 423 & 792 & 2.28 & -- \\
헝가리안 & 452 & 511 & 378 & 447 & 416 & 727 & 2.00 & $-66\pm19$ (15/18) \\
CEM (오류 조건 학습) & 402 & 474 & 496 & 457 & 416 & 675 & 2.28 & $-57\pm22$ (15/18) \\
CEM (오류 없이 학습) & 369 & 459 & 374 & 400 & 415 & 622 & 2.33 & $-113\pm19$ (18/18) \\
\midrule
\multicolumn{9}{l}{\emph{오류 없음, 시드 48--50 (9개)}} \\
팀 기본 WES & 541 & 465 & 499 & 502 & 418 & 715 & 2.11 & -- \\
헝가리안 & 452 & 471 & 425 & 449 & 412 & 672 & 2.22 & $-53\pm32$ (5/9) \\
CEM (오류 조건 학습) & 352 & 416 & 337 & 368 & 405 & 545 & 2.00 & $-133\pm31$ (9/9) \\
CEM (오류 없이 학습) & 354 & 419 & 437 & 403 & 413 & 599 & 2.11 & $-98\pm31$ (8/9) \\
\bottomrule
\end{tabular}
\end{table*}

\begin{itemize}
  \item 두 학습 배정 모두 팀 기본 WES보다 주문 리드타임과 p90이 짧았다. 오류 조건에서 오류 없이 학습한 CEM은 113\,초(22\%) 짧았고 18개 시나리오 모두에서 나았다. 오류 조건에서 학습한 CEM은 57\,초(11\%) 짧았고 18개 중 15개에서 나았다.
  \item 03-09와 03-22 절반에서는 두 학습 배정이 12개 시나리오 모두에서 팀 기본 WES보다 나았다.
  \item 03-22 전체(과부하)는 주문이 13개뿐이라, 큰 주문 한두 개를 끝내느냐에 따라 주문 리드타임이 약 250\,초와 500\,초 사이에서 갈린다. 오류 조건에서 학습한 CEM은 이 부하에서 시드 6개 중 3개에서 큰 주문을 끝내지 못해 팀 기본 WES보다 나빴다(평균 $+29$\,초).
  \item 헝가리안은 경량 시뮬레이터 예측(오류 조건 $-23$\,초)보다 Webots에서 훨씬 강했다($-66$\,초). 오류 없이 학습한 CEM은 오류 조건에서 헝가리안보다 47$\pm$11\,초 짧았지만(14/18), 오류 조건에서 학습한 CEM은 헝가리안과 차이가 없었다($+10\pm27$\,초).
  \item 경량 시뮬레이터에서 본 \emph{오류가 있을 때 학습의 이득이 더 크다}는 Webots에서 확인되지 않았다. 학습 배정의 개선 폭은 오류 유무보다 어떤 가중치를 쓰느냐에 더 크게 좌우되었다.
  \item 상자 단위 리드타임은 방법 간 거의 같았다(405--423\,초). 개선은 처리량이 아니라 \emph{주문을 끝내는 순서}에서 나온다.
\end{itemize}

\subsection{제거 실험}

\begin{table}[t]
\centering
\caption{오류 없음, 경량 시뮬레이터 시험 결과 (45 에피소드 평균). 시간은 초.} % logs/train_nofault_m100*_report.json
\label{tab:nofault}
\small
\begin{tabular}{lrrrr}
\toprule
방법 & 주문 & 상자 & p90 & 마감 놓침 \\
\midrule
팀 기본 WES & 495 & 416 & 712 & 2.22 \\
가까운 로봇 규칙 & 491 & 410 & 678 & 2.33 \\
헝가리안 & 495 & 410 & 677 & 2.3 \\ % scratchpad hung_screen.py, 2026-11-26
CEM & 473 & 412 & 613 & 2.22 \\
\midrule
CMA-ES & 481 & 419 & 620 & 1.89 \\
CEM + 도착 예상 시간 & 477 & 415 & 601 & 1.89 \\
CMA-ES + 도착 예상 시간 & 484 & 419 & 660 & 1.89 \\
항상 확정 + CEM & 490 & 421 & 626 & 1.9 \\
조건부 확정 + CEM$^\dagger$ & 465 & 411 & 570 & 2.1 \\
\bottomrule
\end{tabular}

\vspace{2pt}
\footnotesize $^\dagger$ 시험 중 예약이 한 번도 일어나지 않아, 확정이 없는 정책과 결과가 같다. 차이는 학습 결과의 흔들림이다.
\end{table}

\textbf{다음 작업 확정.}
\begin{itemize}
  \item 항상 확정은 확정 없는 CEM보다 주문 리드타임이 17\,초 나빴다(473$\to$490). 놀고 있는 로봇이 할 수 있는 일을 바쁜 로봇이 붙잡아 두기 때문이다.
  \item 조건부 확정으로 학습하면 예약이 시험 에피소드에서 0회였다. 같은 가중치로 확정을 끄고 돌려도 결과가 똑같았다.
  \item 확정이 일어날 기회가 많은 큰 창고 + 작업자 오류에서도, 학습 시드 3개의 예약 수는 에피소드당 0, 0.07, 0.69회였고 결과는 확정 없는 CEM과 같았다(477$\pm$1 대 477$\pm$4\,초).
  \item 학습하지 않은 규칙에 조건부 확정을 붙이면 큰 창고에서 에피소드당 15.8회 예약했고, p90이 701에서 793\,초로 나빠졌다.
\end{itemize}
즉 학습은 확정을 쓰지 않는 쪽으로 수렴했다.
이 창고에서는 놀고 있는 로봇이 대부분 남아 있어서(가동률 기본 0.59, 큰 창고 0.48), 곧 끝나는 로봇을 기다리는 것보다 놀고 있는 로봇을 바로 보내는 것이 낫다.

\textbf{탐색 방법과 특징.}
CMA-ES나 도착 예상 시간 특징을 써도 CEM과 비슷하거나 나빴다.
학습된 거리 가중치는 0 근처(CEM: $-0.01$) 또는 양수였고,
놀고 있는 로봇이 있는데 일하는 로봇이 작업을 이긴 경우는 시나리오당 0--3회였다.
따라서 개선은 \emph{어느 로봇에게} 줄지가 아니라 \emph{어떤 작업을 먼저} 할지(작업 종류, 대기 시간, 마감 여유에 대한 가중치)에서 나온다.

% ---------------------------------------------------------------
\section{한계}
\begin{itemize}
  \item \textbf{시뮬레이션만 사용한다.} 실제 로봇 시연은 \todo{축소 창고 시연 추가 예정}.
  \item \textbf{오류율은 합성 가정이다.} 실측 오류율 자료가 없어 오류 없음과 기본값 두 단계만 비교했다.
  \item \textbf{학습된 가중치는 하나로 정해지지 않는다.} 작업 종류 특징은 차이만 의미가 있어서, 서로 다른 가중치 조합이 비슷한 성능을 낸다. 가중치 값 하나하나를 해석하는 데는 주의가 필요하다.
  \item \textbf{오래된 작업을 미룰 수 있다.} 평균 비용으로 학습하면 age 가중치가 음수가 되어, 일부 주문이 오래 남을 수 있다. 오류 없는 조건의 Webots 예비 실험에서 학습한 가중치 하나(CMA-ES + 도착 예상 시간)는 과부하 조건에서 주문 하나를 끝내지 못했다.
  \item \textbf{Webots 시험 규모가 작다.} 시나리오당 주문이 13--35개뿐이라, 끝내지 못한 주문 하나가 평균을 수십--수백 초 움직인다. 시나리오마다 1회만 실행했다.
  \item \textbf{학습된 가중치에 따라 결과가 다르다.} Webots에서는 오류 없이 학습한 가중치가 오류 조건에서 학습한 가중치보다 나았고, 이는 경량 시뮬레이터의 예측과 다르다. Webots에서 검증한 가중치는 학습 시드 하나(r0)뿐이다. \todo{r1, r2 Webots 검증}
  \item \textbf{확정의 결론은 이 창고의 가동률에 한정된다.} 로봇이 거의 쉬지 않는 더 높은 부하에서는 결과가 다를 수 있다.
  \item \textbf{경량 시뮬레이터 보정 시점.} 보정은 Webots에 비켜서기 동작을 넣기 전에 했다.
\end{itemize}

% ---------------------------------------------------------------
\section{결론}
작업자가 기록 없이 재고를 옮기는 물류창고에서, 작업-로봇 짝 점수의 가중치를 경량 시뮬레이터에서 CEM으로 학습했다.
학습한 배정은 Webots 시험 시나리오에서 팀 기본 WES보다 주문 리드타임을 11--22\% 줄였으나, 개선 폭은 학습된 가중치에 따라 달랐고 과부하에서는 큰 주문을 놓치는 경우가 있었다.
개선은 어느 로봇에게 줄지가 아니라 어떤 작업을 먼저 할지에서 나왔고, 일하는 로봇에게 다음 작업을 예약하는 확정은 도움이 되지 않았다.
\todo{향후 과제: 확인 작업 결정의 학습, 축소 창고 실물 시연}

% ---------------------------------------------------------------
% 참고문헌: docs/related_work_matrix.md에서 확인한 것과 널리 알려진 고전 문헌만 넣었다.
% 쪽수, 권호는 Overleaf로 옮기기 전에 원문으로 다시 확인할 것.
\begin{thebibliography}{99}
\bibitem{akkerman2025} F. Akkerman, D. Prak, M. Mes, ``Dynamic reordering and inspection for the multi-item inventory record inaccuracy problem,'' \emph{European Journal of Operational Research}, vol.~321, no.~2, 2025. doi:10.1016/j.ejor.2024.09.033
\bibitem{assis2023} de Assis, de Paula Ferreira, Ouhimmou, ``Order Picking Dataset from a Warehouse of a Footwear Manufacturing Company,'' Mendeley Data, V1. doi:10.17632/pf2w725pw3.1 \todo{저자 이름 전체, 연도 확인}
\bibitem{choudhury2020} S. Choudhury \emph{et al.}, ``Dynamic multi-robot task allocation under uncertainty and temporal constraints,'' in \emph{Robotics: Science and Systems (RSS)}, 2020.
\bibitem{dehoratius2008} N. DeHoratius, A.~J. Mersereau, L. Schrage, ``Retail inventory management when records are inaccurate,'' \emph{Manufacturing \& Service Operations Management}, vol.~10, no.~2, 2008.
\bibitem{dhanaraj2024} N. Dhanaraj \emph{et al.}, ``Multi-robot task allocation under uncertainty via hindsight optimization,'' in \emph{IEEE ICRA}, 2024. doi:10.1109/ICRA57147.2024.10611370
\bibitem{gerkey2004} B.~P. Gerkey, M.~J. Matari\'{c}, ``A formal analysis and taxonomy of task allocation in multi-robot systems,'' \emph{The International Journal of Robotics Research}, vol.~23, no.~9, 2004.
\bibitem{hansen2001} N. Hansen, A. Ostermeier, ``Completely derandomized self-adaptation in evolution strategies,'' \emph{Evolutionary Computation}, vol.~9, no.~2, 2001.
\bibitem{kok2007} A.~G. K\"{o}k, K.~H. Shang, ``Inspection and replenishment policies for systems with inventory record inaccuracy,'' \emph{Manufacturing \& Service Operations Management}, vol.~9, no.~2, 2007.
\bibitem{kuhn1955} H.~W. Kuhn, ``The Hungarian method for the assignment problem,'' \emph{Naval Research Logistics Quarterly}, vol.~2, 1955.
\bibitem{leet2023} C. Leet \emph{et al.}, ``Task assignment, scheduling, and motion planning for automated warehouses for million product workloads,'' in \emph{IEEE/RSJ IROS}, 2023.
\bibitem{rubinstein1999} R.~Y. Rubinstein, ``The cross-entropy method for combinatorial and continuous optimization,'' \emph{Methodology and Computing in Applied Probability}, vol.~1, 1999.
\bibitem{wurman2008} P.~R. Wurman, R. D'Andrea, M. Mountz, ``Coordinating hundreds of cooperative, autonomous vehicles in warehouses,'' \emph{AI Magazine}, vol.~29, no.~1, 2008.
\end{thebibliography}

\end{document}
